\chapter[Аштавакра]{Звільнення від ототожнення (Аштавакра)}

[9.1] Протилежні сили --- зроблене й те, що ще слід зробити, --- коли вони припиняються? І для кого? Розмірковуючи над цим, відпусти бажання. Дозволь речам піти. Будь байдужим до мирського.

[9.2] Рідкісний і благословенний той, у кому згасла прив'язаність до насолод і мирського знання після споглядання егоїстичних шляхів людства.

[9.3] Побачивши в мирському джерело потрійного страждання, мудрець залишається незворушним. Він лише споглядає світ --- позбавлений самостійної сутності, плинний і непостійний.

[9.4] Чи був хоч якийсь час, коли людство могло існувати без протилежностей? Забудь про протилежності. Приймай усе, що приходить. Це --- досконалість.

[9.5] Як можна не досягти справжньої незворушності та спокою, помітивши розбіжності в поглядах великих провидців, святих і йогів?

[9.6] Хіба не справжній учитель той, хто завдяки відмові від ототожнення з мирським, спокою та розуму бачить свою істинну природу?

[9.7] У незліченних формах Всесвіту побач сутнісну порожнечу. І ти миттєво знайдеш звільнення на шляху до Я.

[9.8] Бажання породжує оману. Зречися. Зречися бажань і оманливого мирського світу. Тоді ти зможеш пережити, ким є насправді.
